% Présentation du Projet I : enquête et note technique IFC
% Beamer 16:9, style épuré
\documentclass[aspectratio=169,11pt]{beamer}

\usepackage[T1]{fontenc}
\usepackage[utf8]{inputenc}
\usepackage[french]{babel}
\usepackage{lmodern}
\usepackage{booktabs}
\usepackage{colortbl}
\usepackage{tikz}
\usepackage{amsmath,amssymb}

% ---------------------------------------------------------------- couleurs
\definecolor{orangeC}{HTML}{C55A11}
\definecolor{orangeM}{HTML}{F4B183}
\definecolor{orangeL}{HTML}{FDEADA}
\definecolor{encre}{HTML}{3B3B3B}
\definecolor{creme}{HTML}{FFFDFA}

% ---------------------------------------------------------------- thème
\usetheme{default}
\setbeamertemplate{navigation symbols}{}
\setbeamercolor{normal text}{fg=encre,bg=creme}
\setbeamercolor{structure}{fg=orangeC}
\setbeamercolor{frametitle}{fg=white,bg=orangeC}
\setbeamerfont{frametitle}{series=\bfseries,size=\large}
\setbeamertemplate{frametitle}{%
  \vspace*{0.25cm}%
  \begin{beamercolorbox}[wd=\paperwidth,ht=0.85cm,dp=0.3cm,leftskip=0.8cm]{frametitle}
    \usebeamerfont{frametitle}\insertframetitle
  \end{beamercolorbox}}
\setbeamertemplate{itemize item}{\textcolor{orangeC}{\raisebox{0.08em}{\rule{0.5em}{0.5em}}}}
\setbeamertemplate{itemize subitem}{\textcolor{orangeM}{\raisebox{0.08em}{\rule{0.4em}{0.4em}}}}
\setbeamertemplate{footline}{%
  \begin{beamercolorbox}[wd=\paperwidth,ht=0.5cm,dp=0.2cm,leftskip=0.8cm,rightskip=0.8cm]{footline}
    \color{encre}\scriptsize Groupe 1 \hfill Cours d'Assurance Vie, ENSAE ISE 3
    \hfill \insertframenumber\,/\,\inserttotalframenumber
  \end{beamercolorbox}}
\setbeamercolor{block title}{fg=white,bg=orangeC}
\setbeamercolor{block body}{fg=encre,bg=orangeL}
\setbeamertemplate{blocks}[rounded]

% grand chiffre façon vignette
\newcommand{\stat}[2]{%
  \begin{beamercolorbox}[wd=\linewidth,ht=1.5cm,dp=0.8cm,center,rounded=true]{statbox}%
    \centering{\LARGE\bfseries\textcolor{orangeC}{#1}}\\[0.15cm]
    {\small\textcolor{encre}{#2}}%
  \end{beamercolorbox}}
\setbeamercolor{statbox}{bg=orangeL}

% pages de section
\AtBeginSection[]{%
  {\setbeamertemplate{footline}{}%
  \begin{frame}
    \begin{tikzpicture}[remember picture,overlay]
      \fill[orangeC] (current page.south west) rectangle (current page.north east);
      \node[white,font=\Huge\bfseries,align=center] at (current page.center)
        {\insertsectionhead};
      \node[orangeL,font=\large] at ([yshift=-1.6cm]current page.center)
        {\insertsectionnumber\ /\ \totalvalue};
    \end{tikzpicture}
  \end{frame}}}
\newcommand{\totalvalue}{5}

\title{Enquête auprès des institutions d'assurance}
\author{Groupe 1}
\date{Juillet 2026}

\begin{document}

% ---------------------------------------------------------------- titre
{\setbeamertemplate{footline}{}%
\begin{frame}
  \begin{tikzpicture}[remember picture,overlay]
    \fill[orangeC] (current page.north west) rectangle
      ([yshift=-2.9cm]current page.north east);
    \fill[orangeM] ([yshift=-2.9cm]current page.north west) rectangle
      ([yshift=-3.05cm]current page.north east);
    \node[white,font=\bfseries\Large,align=center] at
      ([yshift=-1.15cm]current page.north) {ENQUÊTE AUPRÈS DES INSTITUTIONS D'ASSURANCE};
    \node[white,font=\normalsize,align=center] at
      ([yshift=-2.15cm]current page.north)
      {Note technique du produit le plus commercialisé :\ l'Indemnité de Fin de Carrière (IFC / IFC+)};
  \end{tikzpicture}
  \vspace{2.4cm}
  \begin{columns}[T]
    \begin{column}{0.52\textwidth}
      {\color{orangeC}\itshape Présenté par le Groupe 1 :}\\[0.2cm]
      Gado Giovanni Jocelyn\\
      Alagbe Abdou Hamid\\
      Sié Rachid Traoré\\
      Cheikh Sadibou Ngom
    \end{column}
    \begin{column}{0.42\textwidth}
      {\color{orangeC}\itshape Sous la supervision de :}\\[0.2cm]
      Mme Fatou Diop\\[0.5cm]
      {\small Cours d'Assurance Vie\\ ENSAE Pierre Ndiaye, ISE 3\\ Juillet 2026}
    \end{column}
  \end{columns}
\end{frame}}

% ---------------------------------------------------------------- plan
\begin{frame}{Plan de la présentation}
  \vspace{0.3cm}
  \begin{enumerate}\itemsep0.35cm
    \item L'enquête et ses résultats
    \item Le produit IFC / IFC+
    \item Méthodologie d'évaluation
    \item Application : SENIA S.A.
    \item Le simulateur Excel
  \end{enumerate}
\end{frame}

% ================================================================ SECTION 1
\section{L'enquête et ses résultats}

\begin{frame}{Objectif et démarche}
  \vspace{0.2cm}
  \textbf{Énoncé du projet :} recueillir les produits commercialisés par les
  institutions d'assurance et produire la note technique du produit le plus
  commercialisé.
  \vspace{0.3cm}
  \begin{itemize}\itemsep0.25cm
    \item 8 compagnies d'assurance vie de la place de Dakar interrogées
          (actuaire ou directeur technique) ;
    \item questionnaire en 4 volets : identification, portefeuille de produits,
          caractéristiques techniques, fiche de collecte des données du personnel ;
    \item trois livrables : la note technique, le simulateur Excel et le questionnaire.
  \end{itemize}
\end{frame}

\begin{frame}{Le produit le plus commercialisé}
  \vspace{0.2cm}
  \begin{columns}[T]
    \begin{column}{0.58\textwidth}
      \small
      \rowcolors{2}{white}{orangeL}
      \begin{tabular}{lcc}
        \toprule
        \rowcolor{orangeM} Produit & Cité n\textsuperscript{o}1 & Part moyenne \\
        \midrule
        IFC / IFC+ & 5/8 & 34\,\% \\
        Décès emprunteur & 2/8 & 22\,\% \\
        Épargne / capitalisation & 1/8 & 18\,\% \\
        Temporaire décès & 0/8 & 11\,\% \\
        Retraite complémentaire & 0/8 & 9\,\% \\
        Autres & 0/8 & 6\,\% \\
        \bottomrule
      \end{tabular}
    \end{column}
    \begin{column}{0.38\textwidth}
      \stat{IFC / IFC+}{produit vie n\textsuperscript{o}1 du marché}
      \vspace{0.25cm}
      {\small Un marché porté par l'obligation conventionnelle (CCNI) et
      comptable (SYSCOHADA révisé, IAS 19) de provisionner le passif social.}
    \end{column}
  \end{columns}
\end{frame}

% ================================================================ SECTION 2
\section{Le produit IFC / IFC+}

\begin{frame}{Garanties du contrat}
  \vspace{0.2cm}
  \begin{columns}[T]
    \begin{column}{0.48\textwidth}
      \begin{block}{Volet IFC}
        Constitution et gestion du fonds destiné au paiement des indemnités de
        fin de carrière dues au départ à la retraite, avec évaluation
        actuarielle du passif social.
      \end{block}
    \end{column}
    \begin{column}{0.48\textwidth}
      \begin{block}{Volet IFC Plus}
        Prise en charge de l'indemnité de décès due par l'employeur en cas de
        décès d'un salarié en activité, sans limitation du nombre de décès
        dans l'année.
      \end{block}
    \end{column}
  \end{columns}
  \vspace{0.4cm}
  \textbf{Assurés :} l'ensemble du personnel salarié, adhésion collective,
  sans sélection médicale.
\end{frame}

\begin{frame}{Cadre réglementaire et bases techniques}
  \vspace{0.2cm}
  \begin{columns}[T]
    \begin{column}{0.46\textwidth}
      \textcolor{orangeC}{\textbf{Cadre juridique}}
      \begin{itemize}\itemsep0.15cm
        \item CCNI du Sénégal : barème des indemnités ;
        \item code du travail et IPRES : retraite à 60 ans ;
        \item code CIMA (art. 338) : bases techniques ;
        \item SYSCOHADA révisé / IAS 19.
      \end{itemize}
    \end{column}
    \begin{column}{0.50\textwidth}
      \small
      \rowcolors{2}{orangeL}{white}
      \begin{tabular}{lc}
        \toprule
        \rowcolor{orangeM} Paramètre & Valeur \\
        \midrule
        Table de mortalité & CIMA H \\
        Taux technique & 3,5\,\% \\
        Évolution des salaires & 2,0\,\% \\
        Turn-over & 2,0\,\% \\
        Licenciement annuel & 1,0\,\% \\
        Frais de gestion & 2,0\,\% \\
        \bottomrule
      \end{tabular}
    \end{column}
  \end{columns}
\end{frame}

\begin{frame}{Barème CCNI des droits}
  \vspace{0.2cm}
  Droit exprimé en \% du salaire global mensuel moyen, par année de présence :
  \vspace{0.3cm}
  \begin{center}
    \rowcolors{2}{white}{orangeL}
    \begin{tabular}{lc}
      \toprule
      \rowcolor{orangeM} Tranche d'ancienneté & Droit par année \\
      \midrule
      Années 1 à 5 & 20\,\% \\
      Années 6 à 10 & 30\,\% \\
      Années 11 à 15 & 40\,\% \\
      Années 16 à 20 & 50\,\% \\
      Au-delà de 20 ans & 60\,\% \\
      \bottomrule
    \end{tabular}
  \end{center}
  \vspace{0.2cm}
  \small Cumul $B(a)$ obtenu par sommation, avec interpolation linéaire pour
  les anciennetés non entières.
\end{frame}

% ================================================================ SECTION 3
\section{Méthodologie d'évaluation}

\begin{frame}{Méthode prospective}
  \vspace{0.2cm}
  Droits projetés au terme, probabilisés, actualisés, puis répartis au prorata
  de l'ancienneté.
  \vspace{0.25cm}
  \begin{block}{Salaire final et droits au terme}
    \vspace{-0.25cm}
    $$ S_f = \frac{SA}{12}\,(1+\tau_s)^{\max(60-x,\,0)}
       \qquad\qquad \text{Droits} = B(N_r)\times S_f $$
  \end{block}
  \begin{block}{Valeur actuelle probable des prestations futures}
    \vspace{-0.25cm}
    $$ \mathrm{VAP} = \text{Droits}\times\frac{l_{60}}{l_x}\times v^{\,n+1/2} $$
  \end{block}
  \small $x$ : âge ; $a$ : ancienneté ; $n$ : durée résiduelle ;
  $N_r = a + n$ : ancienneté au terme ; $v=(1+i)^{-1}$, $i=3{,}5\,\%$.
\end{frame}

\begin{frame}{De l'engagement au passif comptable}
  \vspace{0.2cm}
  \begin{block}{Dette actuarielle et charge normale}
    \vspace{-0.25cm}
    $$ \mathrm{DA} = \mathrm{VAP}\times\frac{a}{N_r}
       \qquad\qquad \mathrm{CN} = \frac{\mathrm{VAP}}{N_r} $$
  \end{block}
  \begin{block}{Passif de l'indemnité de licenciement}
    \vspace{-0.25cm}
    $$ \mathrm{Passif\ IL} = \mathrm{IL}\times\tau_\ell\times\frac{1-u^{\,n}}{1-u},
       \qquad u=\frac{1-\tau_r}{1+i} $$
  \end{block}
  \begin{block}{Tarification IFC+ (décès)}
    \vspace{-0.25cm}
    $$ \mathrm{PP} = \sum_j C_j\,q_{x_j}
       \qquad\qquad \mathrm{PC} = \frac{\mathrm{PP}}{1-g} $$
  \end{block}
\end{frame}

% ================================================================ SECTION 4
\section{Application : SENIA S.A.}

\begin{frame}{Le groupe assuré au 31/12/2025}
  \vspace{0.3cm}
  \begin{columns}[T]
    \begin{column}{0.31\textwidth}\stat{50}{salariés, 5 catégories}\end{column}
    \begin{column}{0.31\textwidth}\stat{344,4 M}{masse salariale (FCFA)}\end{column}
    \begin{column}{0.31\textwidth}\stat{42,2 ans}{âge moyen}\end{column}
  \end{columns}
  \vspace{0.35cm}
  \begin{columns}[T]
    \begin{column}{0.31\textwidth}\stat{10,2 ans}{ancienneté moyenne}\end{column}
    \begin{column}{0.31\textwidth}\stat{17,8 ans}{durée résiduelle moyenne}\end{column}
    \begin{column}{0.31\textwidth}\stat{60 ans}{âge de départ (IPRES)}\end{column}
  \end{columns}
  \vspace{0.35cm}
  \small Base recueillie auprès de l'entreprise à l'aide de la fiche de
  collecte du questionnaire (volet D).
\end{frame}

\begin{frame}{Résultats de l'évaluation IFC}
  \vspace{0.2cm}
  \begin{columns}[T]
    \begin{column}{0.56\textwidth}
      \rowcolors{2}{white}{orangeL}
      \begin{tabular}{lr}
        \toprule
        \rowcolor{orangeM} Rubrique & FCFA \\
        \midrule
        Engagement global & 230\,047\,411 \\
        Dette actuarielle & 111\,103\,527 \\
        Charge normale & 8\,223\,109 \\
        Passif indemnité de licenciement & 7\,831\,464 \\
        \bottomrule
      \end{tabular}
    \end{column}
    \begin{column}{0.40\textwidth}
      \stat{230 M FCFA}{d'engagement global}
      \vspace{0.2cm}
      {\small La dette actuarielle représente 48\,\% de l'engagement,
      en cohérence avec le prorata d'ancienneté (10,2 / 28,1 ans).}
    \end{column}
  \end{columns}
\end{frame}

\begin{frame}{Modalités de financement}
  \vspace{0.2cm}
  \rowcolors{2}{orangeL}{white}
  \begin{tabular}{lr}
    \toprule
    \rowcolor{orangeM} Option & FCFA \\
    \midrule
    Option 1 : prime unique (engagement + passif IL) & 237\,878\,875 \\
    Option 2 : dette actuarielle à la souscription\ldots & 111\,103\,527 \\
    \quad\ldots puis charge normale annuelle & 8\,223\,109 \\
    Option 3 : amortissement de la dette sur 5 ans & 22\,220\,705 \\
    \quad prime annuelle sur 5 ans & 30\,443\,814 \\
    \bottomrule
  \end{tabular}
  \vspace{0.3cm}

  \small L'option 3 (taux de prime uniforme de 8,84\,\% de la masse salariale)
  permet la déduction fiscale intégrale de la prime ; réactualisation de
  l'étude au minimum tous les trois ans.
\end{frame}

\begin{frame}{Volet IFC+ et sensibilité}
  \vspace{0.2cm}
  \begin{columns}[T]
    \begin{column}{0.44\textwidth}
      \textcolor{orangeC}{\textbf{Tarification IFC+ (an 1)}}\\[0.15cm]
      \small
      \rowcolors{2}{white}{orangeL}
      \begin{tabular}{lr}
        \toprule
        \rowcolor{orangeM} Rubrique & Valeur \\
        \midrule
        Capitaux garantis & 161\,949\,480 \\
        Prime commerciale & 1\,203\,354 \\
        Taux / capitaux & 0,74\,\% \\
        Taux / masse salariale & 0,35\,\% \\
        \bottomrule
      \end{tabular}
    \end{column}
    \begin{column}{0.52\textwidth}
      \textcolor{orangeC}{\textbf{Sensibilité de l'engagement}}\\[0.15cm]
      \small
      \rowcolors{2}{white}{orangeL}
      \begin{tabular}{lr}
        \toprule
        \rowcolor{orangeM} Scénario & Engagement (FCFA) \\
        \midrule
        Central ($i=3{,}5\,\%$, $\tau_s=2\,\%$) & 230\,047\,411 \\
        $i=3{,}0\,\%$ & 249\,478\,037 \\
        $i=4{,}0\,\%$ & 212\,687\,421 \\
        $\tau_s=3\,\%$ & 271\,442\,252 \\
        $\tau_s=1\,\%$ & 196\,416\,468 \\
        \bottomrule
      \end{tabular}
    \end{column}
  \end{columns}
\end{frame}

% ================================================================ SECTION 5
\section{Le simulateur Excel}

\begin{frame}{Organisation du classeur}
  \vspace{0.2cm}
  \rowcolors{2}{orangeL}{white}
  \begin{tabular}{ll}
    \toprule
    \rowcolor{orangeM} Feuille & Contenu \\
    \midrule
    Hypotheses & paramètres de calcul modifiables (dates, taux, frais) \\
    BD & base de données du personnel (50 salariés) \\
    Baremes & barème CCNI des droits \\
    Table CIMA H & table de mortalité et commutations \\
    Calculs & évaluation individuelle : droits, VAP, dette, charge, passif \\
    Resultats & synthèse : statistiques, résultats IFC, financement, IFC+ \\
    \bottomrule
  \end{tabular}
  \vspace{0.3cm}

  \small Formules vivantes de bout en bout : toute modification d'une hypothèse
  ou de la base se propage jusqu'à la synthèse ; le simulateur est réutilisable
  pour toute autre entreprise.
\end{frame}

\begin{frame}{Conclusion}
  \vspace{0.3cm}
  \begin{itemize}\itemsep0.3cm
    \item L'IFC est le produit vie le plus commercialisé de la place, porté par
          les obligations conventionnelles et comptables des entreprises ;
    \item la méthode prospective fournit un passif social complet : engagement
          global, dette actuarielle, charge normale et passif de licenciement ;
    \item trois modalités de financement adaptées à la situation de l'entreprise ;
    \item un simulateur opérationnel, alimenté par le questionnaire de collecte.
  \end{itemize}
\end{frame}

{\setbeamertemplate{footline}{}%
\begin{frame}
  \begin{tikzpicture}[remember picture,overlay]
    \fill[orangeC] (current page.south west) rectangle (current page.north east);
    \node[white,font=\Huge\bfseries] at ([yshift=0.6cm]current page.center)
      {Merci de votre attention};
    \node[orangeL,font=\large] at ([yshift=-0.9cm]current page.center)
      {Groupe 1, Cours d'Assurance Vie, ENSAE ISE 3};
  \end{tikzpicture}
\end{frame}}

\end{document}
